\section{Paired Warm-Start Evaluation}\label{sec:problem-setting}

For each independently prepared student, a validation-selected $S_0$ initializes CE and CE+KD with identical weights and fresh optimizer states. Continuation data, order, optimizer, scheduler, and budget are matched. Each arm selects its checkpoint independently by validation NLL; test reporting follows selection (Figure~\ref{fig:paired_protocol}).

Let $p_{\mathrm{teach},i}^{(\tau)}$ and $p_{\mathrm{stud},i}^{(\tau)}$ denote next-token distributions at temperature $\tau$. Over valid target positions $\mathcal I$ after the causal shift, the objective is
\begin{align}
\mathcal L&=\mathcal L_{\mathrm{CE}}+\lambda\mathcal L_{\mathrm{distill}},\label{eq:kd-objective}\\
\mathcal L_{\mathrm{distill}}&=\frac{\tau^2}{|\mathcal I|}\sum_{i\in\mathcal I}
\mathrm{KL}\!\left(p_{\mathrm{teach},i}^{(\tau)}\|p_{\mathrm{stud},i}^{(\tau)}\right).\nonumber
\end{align}
Both losses average over target tokens. Primary KD fixes $\tau=2$ and $\lambda=5$; CE uses $\lambda=0$. The operational paired endpoint is
\begin{equation}
\mathrm{Gain}_{\mathrm{KD}}=\mathrm{NLL}_{\mathrm{CE}}^{\mathrm{test}}-\mathrm{NLL}_{\mathrm{KD}}^{\mathrm{test}}.
\label{eq:paired-gain}
\end{equation}
Positive gain favors KD over CE, not necessarily over $S_0$. This is a protocol-specific endpoint, not a theoretical measure of distillability.

We report seeds 42, 123, and 456, their mean paired gain, sample SD, and sign consistency. These vary student preparation, not teacher training or corpus sampling; SD is neither a standard error nor a significance test. Primary selection uses trained continuation checkpoints. A separately labeled diagnostic allows unmodified $S_0$ as step 0 without replacing the primary endpoints.
